\section{Vai trò của chatbot Rasa trong hệ thống}

Trong hệ thống KMA Chatbot tuyển sinh, Rasa đóng vai trò là phân hệ hội thoại chính dùng để xử lý các câu hỏi phổ biến, các luồng hỏi đáp định sẵn và các nghiệp vụ có cấu trúc rõ ràng. Khác với mô hình sinh văn bản tự do, Rasa được thiết kế theo hướng kiểm soát luồng hội thoại, trong đó mỗi câu hỏi của người dùng được phân tích thành ý định, trạng thái hội thoại và hành động phản hồi tương ứng.

Trước khi phân tích kiến trúc Rasa sử dụng NLU và PhoBERT, cần phân biệt với mô hình chatbot dựa trên luật. Chatbot rule-based vận hành chủ yếu dựa vào tập quy tắc if--then, từ khóa và các mẫu câu được định nghĩa trước. Cách tiếp cận này đơn giản, dễ kiểm soát nhưng kém linh hoạt khi người dùng diễn đạt câu hỏi theo nhiều cách khác nhau.

\begin{figure}[H]
    \centering
    \includegraphics[width=0.82\textwidth]{\detokenize{content/source/rasa/trợ lý ảo dựa trên rule.png}}
    \caption{Kiến trúc chatbot dựa trên luật}
    \label{fig:rasa-rule-based-assistant}
\end{figure}

Hình~\ref{fig:rasa-rule-based-assistant} mô tả kiến trúc chatbot dựa trên luật. Câu hỏi của người dùng được đưa qua bộ xử lý ngôn ngữ để tách từ khóa, chuẩn hóa văn bản và so khớp với tập quy tắc. Nếu tìm được quy tắc phù hợp, hệ thống lấy câu trả lời từ cơ sở dữ liệu hoặc bộ tạo phản hồi. Ưu điểm của kiến trúc này là dễ triển khai và dễ kiểm soát nội dung. Tuy nhiên, khi câu hỏi có nhiều cách diễn đạt khác nhau hoặc chứa ngữ cảnh phức tạp, chatbot rule-based dễ thất bại do phụ thuộc mạnh vào từ khóa và luật được viết thủ công.

Khác với cách tiếp cận rule-based, Rasa kết hợp NLU, intent classification, entity extraction, tracker, policy và action để xử lý hội thoại linh hoạt hơn. Khi tích hợp PhoBERTIntentClassifier, hệ thống không chỉ so khớp từ khóa mà còn học biểu diễn ngữ nghĩa của câu hỏi tiếng Việt, từ đó phân biệt tốt hơn các intent gần nghĩa như hỏi điểm chuẩn, hỏi phương thức xét tuyển, hỏi học phí hoặc hỏi chương trình đào tạo.

\begin{figure}[H]
    \centering
    \includegraphics[width=0.9\textwidth]{\detokenize{content/source/rasa/trợ lý ảo dựa trên NLP.png}}
    \caption{Kiến trúc trợ lý ảo dựa trên NLP và mô hình học máy}
    \label{fig:rasa-nlp-based-assistant}
\end{figure}

Hình~\ref{fig:rasa-nlp-based-assistant} thể hiện hướng tiếp cận chatbot dựa trên NLP và học máy. Câu hỏi của người dùng được phân tích để nhận diện ý định, trích xuất thông tin quan trọng và chuyển thành dữ liệu có cấu trúc cho bộ quản lý hội thoại.

Luồng chat thường của hệ thống bắt đầu từ giao diện hội thoại tại route \texttt{/} hoặc \texttt{/home\_chat}. Người dùng nhập câu hỏi vào giao diện, frontend gửi nội dung về backend chatbot, sau đó backend phối hợp với Rasa để xác định intent, entity, rule/story và action phù hợp. Phạm vi nghiệp vụ chính gồm tư vấn phương thức tuyển sinh, học phí, giấy báo trúng tuyển, hồ sơ nhập học và thông tin ngành đào tạo.

Trong kiến trúc tổng thể, Rasa không trực tiếp thay thế toàn bộ hệ thống hỏi đáp, mà đảm nhiệm phần hiểu ý định và điều phối hội thoại. Khi người dùng hỏi các câu thuộc nhóm đã được định nghĩa sẵn, Rasa có thể phản hồi nhanh bằng template trong domain. Khi câu hỏi cần xử lý logic hoặc không đủ chắc chắn, Rasa chuyển sang custom action hoặc fallback để tránh trả lời sai.

Vì vậy, phân hệ Rasa được thiết kế theo nguyên tắc: câu hỏi đơn giản và phổ biến được xử lý bằng intent, rule và response template; câu hỏi cần logic được xử lý bằng action server; câu hỏi mơ hồ hoặc ngoài phạm vi được đưa vào fallback hoặc ghi log để cải thiện dữ liệu huấn luyện.

